\begin{abstract}
Airline operations control centers (AOCCs) must allocate limited attention to disrupted flight chains and identify where local recovery remains worthwhile. We ask whether decision-time expected successor take-off delay supports both cross-chain screening and within-chain recovery targeting, and when additional aircraft-linked information changes these choices. Using U.S. airline operating data from 2017 to 2022, we examine chronologically held-out cases each year at prespecified PRE (before predecessor arrival) and TURN (after predecessor arrival but before successor departure) epochs, comparing delay- and consequence-based priorities and Point, Marginal, and Joint state representations. At nominal 10\% review capacity, delay priority displaces 37--53\% of PRE and 40--55\% of TURN positions from the consequence-based reference shortlist, with normalized consequence losses of 0.101--0.184 and 0.107--0.181, respectively. Passenger exposure reduces part of this loss, whereas downstream exposure adds little. Preserving weighted marginals while changing cross-state dependence reassigns only 1.3--3.8\% of shortlist positions and produces screening loss of 0.0006--0.0026, yet exact local-recovery target agreement is 0.581--0.730 and local-recovery loss is 0.085--0.125 within annual consequence-based reference cohorts. By contrast, point compression relative to the joint representation reassigns 84.2--88.7\% of shortlist positions and yields local-recovery losses of 0.896--0.981. These findings show that information sufficiency is decision-specific: representations that preserve screening choices need not preserve recovery targets, and information value cannot be inferred from representational richness alone.
\end{abstract}
